\section{Restatements of the Main Results}\label{sec:formal-statement}
In this section, to make our results more organized and easy to understand, we give the restatements of the main results in \Cref{sec:theory}, where we present the two main theorems.
\begin{enumerate}
  \item The AGM iterates converge to a neighborhood of the manifold (\Cref{thm: convergence});
  \item Moreover, once the iterates enter this neighborhood, their dynamics over $\O(\eta^{-2})$ discrete steps can be accurately tracked by a slow SDE (\Cref{thm:SDE-approximation}).
\end{enumerate}


Recall that in the AGM framework, the transition from $\vtheta_{k}$ to $\vtheta_{k+1}$ is defined as:
\begin{align*}
    \vm_{k+1} &:= \beta_1 \vm_{k} + (1-\beta_1) \nabla \ell_{k}(\vtheta_{k}) \\
    \vv_{k+1} &:= \beta_2 \vv_{k} + (1-\beta_2) V\!\left(\nabla\ell_{k}(\vtheta_{k}) \nabla\ell_{k}(\vtheta_{k})^\top\right) \\
    \vtheta_{k+1} &:= \vtheta_{k} - \eta S(\vv_{k+1}) \vm_{k+1},
\end{align*}
under the following conditions:
\begin{enumerate}
    \item $S: \Rgez^{D} \longrightarrow \mathbb{S}^d_{++}$ is $\rho_s$-smooth, where $\Rgez^{D}$ denotes the subset of $\R^D$ that has non-negative entries, and $\mathbb{S}^d_{++}$ denotes the space of $\R^{d \times d}$ positive definite matrices. 
    \item $S(\vv) \succeq \frac{1}{R_0} \mI$  for some $R_0 > 0$ and any $\vv \in \Rgez^{D}$.
    \item $V: \R^{d\times d} \longrightarrow \R^{D}$ is linear, and $V(\vg\vg^\top) \in \Rgez^{D}$ for all $\vg \in \R^d$.
\end{enumerate}


% \begin{assumption}
%     vfk has been awarded the Tsinghua University Special Scholarship for Outstanding Students.
% \end{assumption}

\subsection{Slow SDE for AGMs}

\begin{theorem}\label{thm:formal-main-result}
Let Assumptions \ref{amp:smooth}, \ref{amp:gradient-bound} and \ref{amp:beta1-bound} be satisfied. Let $\Gamma$ denote a local minimizer manifold, $\eta$ be a sufficiently small learning rate of an AGM, and $1-\beta_2 = \Theta(\eta^2)$. 
Then we have the following conclusions:
\begin{enumerate}
    \item (Convergence to a near–manifold neighborhood) There exists a constant $\epsilon_0 > 0$ such that for any initial point $\vtheta_0$ whose $\ell_2$ distance from $\Gamma$ does not exceed $\epsilon_0$, and any $\delta \in (\eta^{200}, 1)$,\footnotemark[3] with probability at least $1 - \delta$, the following holds for some $K_0 = \mathcal{O}(\frac{1}{\eta} \log \frac{1}{\eta})$:
    \begin{gather*}
        \cL(\vtheta_{K_0}) - \cL^* = \O\left(\eta\log\frac{1}{\eta\delta}\right),\\
        \|\vtheta_{K_0} -\Phi_{\mS_{K_0}}(\vtheta_{K_0})\|_2 = \O\left(\sqrt{\eta\log\frac{1}{\eta\delta}}\right).
    \end{gather*} 
    \footnotetext[3]{The exponent here, along with the exponents related to the $\delta$-goodness in \Cref{def:delta-good}, can be arbitrary large constant, which does not affect the order of following derivations.}
    % where the notation $\Gamma^{\epsilon}$ denotes $\left\{\vtheta : \exists \vtheta' \in \Gamma, \lnormtwo{\vtheta - \vtheta'} \leq \epsilon\right\}$ , i.e. the $\epsilon$-neighborhood of $\Gamma$.
    
    \item (Slow SDE tracks AGM's trajectory in a weak approximation sense) Moreover, let $T > 0$, $K := \lfloor T\,\eta^{-2}\rfloor$, and suppose Assumptions \ref{amp:low-rank-manifold}, \ref{amp:basic} and \ref{amp:S-smooth} hold. Continue running an AGM for $K$ iterations after reaching the final state $(\vtheta_{K_0},\vv_{K_0})$ in conclusion 1 and totally obtain $\{(\vtheta_k, \vv_k)\}_{k=0}^{K_0+K}$. Let $\mX(t) =\left(\vzeta(t),\vv(t)\right)$ be the solution to \eqref{equ:AGMs-SDE} with initial condition
$$\mX(0) =
\mX_0
\in \Gamma \times \mathbb{R}_{\ge 0}^D, \quad \mX_0 := (\Phi_{\mS(\vv_{K_0})}(\vtheta_{K_0}), \vv_{K_0}).
$$
% and we define that $\vtheta^{(s)}:=\vtheta_{sH}$, and $\bar\mX^{(s)} := \left(\Phi_{\mS^{(s)}}(\vtheta^{(s)}),\vv^{(s)}\right)$, where $H$ is $\frac{\alpha}{\eta}$ for some constant $\alpha>0$.

% and define the parameters of AGM as $\bar\mX_{t} := \left(\Phi_{\mS_t}(\vtheta_{t}),\vv_{t}\right)$. For any $C^3$ function $g(\vtheta)$,
% \[
%   \max_{0\le t\le \lfloor\frac{T}{\eta^2}\rfloor}
%   \left|\E\left[g\left(\bar\mX_{t}\right) \middle| \mX_0\right]
%   -\E\left[g\left(\mX( t\eta^2)\right)\middle|\mX(0)=\mX_0\right]\right|
%   =\widetilde \O\bigl(\eta^{0.25}\bigr),
% \]
% where $\widetilde O(\cdot)$ hides logarithmic factors and constants that are independent of~$\eta$ (but possibly depending on $g$).
For any $k \in \{0, 1, \dots, K\}$,
let $\bar\mX_{k} := \left(\Phi_{\mS(\vv_{K_0 + k})}(\vtheta_{K_0+k}),\vv_{K_0+k}\right)$ be the AGM state at step $K_0 + k$ projected onto $\Gamma$.
%For any $k \in \{0,1,\dots,K\}$, we denote by $\bar\mX_{k}$ the AGM state at step $K_0 + k$ projected to $\Gamma$, i.e., $$\bar\mX_{k} := \left(\Phi_{\mS(\vv_{K_0 + k})}(\vtheta_{K_0+k}),\vv_{K_0+k}\right).$$
Then for any $C^3$ function $g(\vtheta)$,
\begin{align*}
\max_{0\le k\le K}
\left|
\E\left[g\left(\bar\mX_{k}\right)\middle|\mX_0\right]
-\E\left[g\left(\mX(k\eta^2)\right)\middle|\mX(0)=\mX_0\right]\right| 
= \widetilde \O\left(\eta^{0.25}\right),
\end{align*}
where $\widetilde \O(\cdot)$ hides logarithmic factors and constants independent of $\eta$ (but possibly depending on $g$).
\end{enumerate}

\end{theorem}

% \begin{theorem}
% Let $T>0$ and suppose \Cref{amp:basic}--\ref{amp:beta1-bound} hold. 
% There exist constant $\epsilon_0,C>0$ such that the following statement holds for all sufficiently small learning rates $\eta$.
% Define step $K_0 := \lfloor C\log(1/\eta)\rfloor$ and $K := \lfloor T\,\eta^{-2}\rfloor$. 
% %Assume the initial parameter $\vtheta_0$ satisfies $\mathrm{dist}(\vtheta_0,\Gamma)\le \epsilon_0$.
% Run an AGM for $K_0+K+1$ iterations and obtain $\{(\vtheta_k, \vv_k)\}_{k=0}^{K_0+K}$, where the initial parameter $\vtheta_0$ 
% is an arbitrary point that satisfies $\mathrm{dist}(\vtheta_0,\Gamma)\le \epsilon_0$.
% Let $\mX(t)=(\vzeta(t),\vv(t))$ be the solution to \eqref{equ:AGMs-SDE} on $t\in[0,T]$ with initial condition 
% $$\mX(0) =
% \mX_0
% \in \Gamma \times \mathbb{R}_{\ge 0}^D, \quad \mX_0 := (\Phi_{\mS(\vv_{0})}(\vtheta_{0}), \vv_{0}).
% $$
% For any $k \in \{0, 1, \dots, K\}$,
% let $\bar\mX_{k} := \left(\Phi_{\mS(\vv_{K_0 + k})}(\vtheta_{K_0+k}),\vv_{K_0+k}\right)$ be the AGM state at step $K_0 + k$ projected onto $\Gamma$.
% %For any $k \in \{0,1,\dots,K\}$, we denote by $\bar\mX_{k}$ the AGM state at step $K_0 + k$ projected to $\Gamma$, i.e., $$\bar\mX_{k} := \left(\Phi_{\mS(\vv_{K_0 + k})}(\vtheta_{K_0+k}),\vv_{K_0+k}\right).$$
% Then for any $C^3$ function $g(\vtheta)$,
% \begin{align*}
% \max_{0\le k\le K}
% \left|
% \E\left[g\left(\bar\mX_{k}\right)\middle|\mX_0\right]
% -\E\left[g\left(\mX(k\eta^2)\right)\middle|\mX(0)=\mX_0\right]\right| 
% = \widetilde \O\left(\eta^{0.25}\right),
% \end{align*}
% where $\widetilde \O(\cdot)$ hides logarithmic factors and constants independent of $\eta$ (but possibly depending on $g$).
% \end{theorem}

Notice that in \Cref{thm:formal-main-result} we give a more explicit statement for our SDE approximation results in \Cref{thm:SDE-approximation}. One can see that \Cref{thm:SDE-approximation} is a direct corollary of \Cref{thm:formal-main-result}.

\begin{proof}[Proof of \Cref{thm:formal-main-result}]
    For the convergence part of \Cref{thm:formal-main-result}, we first construct the working zones near the manifold in \Aref{sec:working-zone} to ensure some properties that are crucial to our analysis such as the $\mu$-PL condition. After that, we give the proof of convergence in \Aref{app:proof-convergence}. We finish our proof in \Aref{app:proof-formal-statement} by combining the previous results in working zone and convergence analysis.
    
    And the proof of the SDE approximation part is finished in \Aref{sec:weak-approximation} through the moment calculation and the subsequent weak approximation analysis in \Aref{sec:proof-sde-approx}.
\end{proof}

\subsection{Convergence Guarantee of AGMs}

In the proof, the first part of \cref{thm:formal-main-result} is done by first proving a convergence result with global $\mu$-PL condition, and then arguing that AGM starting near enough to the manifold will stick to the manifold with high probability. As mentioned in \cref{subsubsec:convergence}, the convergence under $\mu$-PL condition can be seen as a separate technical contribution of our paper, which is stated below.

\begin{definition}[Polyak-Łojasiewicz Condition]
    \label{def:mu-pl}
    For some $\mu, \bar{L} > 0$, we say some function $\cL:\R^d \to \R$ is $(\mu, \bar{L})$-Polyak-Łojasiewicz (abbreviated as $(\mu, \bar{L})$-PL), if and only if for all $\vtheta \in \mathbb{R}^d$ such that $\cL(\vtheta) \leq \bar{L}$:
    $$2\mu(\cL(\vtheta)-\cL^*) \leq \lnormtwo{\nabla \cL(\vtheta)}^2.$$
    When $\bar{L} = +\infty$, we call this condition the $\mu$-Polyak-Łojasiewicz ($\mu$-PL) condition.
\end{definition}


\begin{theorem}[Formal restatement of \Cref{thm: convergence}]\label{thm:formal-convergence-1}
    Let Assumptions \ref{amp:smooth}, \ref{amp:gradient-bound} and \ref{amp:beta1-bound} be satisfied, $\cL$ be a function satisfying the $\mu$-PL condition, and $\eta$ be a sufficiently small learning rate of an AGM. Let $1-\beta_2 = \Theta(\eta^2)$. For any $\delta \in (0,1)$, with probability at least $1 - \delta$, the following holds for some $K = \mathcal{O}(\frac{1}{\eta} \log \frac{1}{\eta})$:
    \begin{gather*}
        \cL(\vtheta_K) - \cL^* = \O\left(\eta\log\frac{1}{\eta\delta}\right),\\
        \|\vtheta_K -\Phi_{\mS_K}(\vtheta_K)\|_2 = \O\left(\sqrt{\eta\log\frac{1}{\eta\delta}}\right).
    \end{gather*} 
\end{theorem}

\paragraph{Remark.} Note that \cref{thm:formal-convergence-1} is different from Part 1 of \cref{thm:formal-main-result} in the sense that \cref{thm:formal-convergence-1} requires the $\mu$-PL condition to hold globally, while Part 1 of \cref{thm:formal-main-result} does not. Actually, the latter requires the iteration to start from some neighborhood of $\Gamma$. Later on, we will find from \cref{lem:working-zone} that $\mu$-PL provably exists in a neighborhood of $\Gamma$, and we prove that an iteration of AGMs starting within that neighborhood stays within that neighborhood with high probability.

There have been many previous works discussing the convergence bound of Adam. For example, \citet{reddi2018convergence} and \citet{dereich2024convergence} give convergence bounds under the convexity condition, \citet{zou2019sufficient}, \citet{shi2021rmsprop} and \citet{zhang2022adam} focus on the cases where learning rates follow a $1/\sqrt{t}$ decay, and the bounds given by \citet{NEURIPS2018_90365351}, \citet{zhang2022adam} and \citet{10.1145/3637528.3671718} do not decrease to $0$ as $\eta \to 0$. Also, most works \citep{defossez2020simple, guo2025unifiedconvergenceanalysisadaptive, iiduka2022theoreticalanalysisadamusing, wang2024convergenceadamnonuniformsmoothness, zhang2024convergence, hong2023highprobabilityconvergenceadam} only establish an upper bound on the average of gradient norms over the time of iteration. In contrast, we directly bound the loss term of the last step to $o(1)$. Going beyond convex loss functions, we establish the bound on $\mu$-PL functions, and we focus on the constant learning rate schedule.
